\section*{Executive Summary}
\noindent\textbf{To: Mr. Timothy Ryan, University Fire Marshal and Director of Fire Protection and Emergency Services}\par
\noindent\textit{Addressee specified by the supplied contest brief. Desk review dated 9 November 2026.}

We studied the evacuation of all 1,329 modelled audience members, including three wheelchair users, using reconstructed circulation, individual response and route choice, finite queues, six Crowd Managers and optional LED guidance/open belts. The main questions are the effects of concert-emergency guidance and interventions; lecture familiarity is represented explicitly but its demographic and response assumptions are not measured at Bailey.

A new manager-message comparison evaluated 100 paired fresh audiences for each of four policies, with the same six posts and final-forty commitment rule. Nearest-time instructions from all six managers averaged 162.4 s (nominal 95\% interval 159.7--165.0 s), with 92/100 clearances by 180 s. Lower-floor-only nearest-time instructions averaged 178.7 s; fixed coordinated destinations averaged 177.2 s. All 400 runs cleared the complete audience. No walking, response, compliance or capacity parameter was retuned. No added LEDs or barriers were used in this comparison.

\textbf{The tested corridor belts and full partitions should not be installed on the basis of this study.} In the earlier 500-run intervention study, partial belts had no clear additional benefit over LEDs (paired change $+0.54$ s; 95\% interval $-0.42$ to $+1.50$ s). They narrow circulation near auditorium-door junctions. Sightline preservation and a graph bypass do not address door swings, post bases, cane detection or crowd contact. Current Cornell EHS guidance requires unobstructed egress circulation. The experimental drawings are not installation plans.

\textbf{Prioritize a venue walkthrough of clear, coordinated staff instructions and supplementary directional signs, keeping circulation unobstructed.} Every proposed LED needs an intelligible local arrow and continuity to a verified exterior route, with particular scrutiny of stage-side access and balcony stairs. The simulation's exit-name messages and staff-priority rule do not validate arrow comprehension, consistency between devices, emergency power or failure behavior. Any physical pilot must be integrated with the venue's existing approved guidance.

The 180 s mean is conditional, not an operational guarantee. The earlier intervention study's slower-response sensitivity averaged 258.8 s with its combined policy, and finer integration changed individual clearances by up to 19.1 s in three checks. Geometry, response, correlated attention, external landing congestion and assistance logistics remain uncertain. No measured Bailey evacuation was available for calibration.

The accompanying manuscript documents these limitations, literature, layout review, tests and preserved experiment evidence. The software review fixes installation-validation and interface defects while retaining the historical results and reproducibility. The next decision requires real measurements and a supervised comprehension/route walkthrough, not further parameter changes to force a faster simulation.
